% GENERATED FILE. Edit canonical lessons, then run npm run generate:books.
% canonical-source-hash: fnv1a64:b1591585469e79d7
% canonical-lessons: SA-C54-iti, SA-C54-iti-clause, SA-R54-iti

\chapter{Thus It Was Said}
\label{ch:sa-thus-it-was-said}
\hlchaptermodality{54}

\begin{chapteropening}
\textbf{By the end of this chapter:} \emph{I can close a quoted clause with \sk{इति} and make a whole Sanskrit sentence the object of a verb of thinking or saying.}

Complete SA-R54-iti and name the two joining mechanisms the track now has, saying which of them subordinates.
\end{chapteropening}

\section[iti]{\sk{इति} (iti) — thus}

\label{lesson:SA-C54-iti}



\begin{quote}
Say both correlative sentences — the one who goes, and where the village is. Then say \emph{cintayati} and what it means.
\end{quote}

\subsection*{You'll want to know: \sk{इति}}
\textbf{\sk{इति}} (\emph{iti}) — \textquotedblleft{}thus, so\textquotedblright{}.

Sanskrit has no quotation marks. It has this word instead, and it does the job better than a pair of commas in the air: \textbf{\sk{इति} stands after the quoted words and closes them.}

\begin{quote}
\textbf{… \sk{इति}} — \emph{\textquotedblleft{}…\textquotedblright{} — that is what was said}
\end{quote}

Whatever comes before \textbf{\sk{इति}} is quoted; whatever comes after it is the frame. Everything the frame needs to say — who said it, who thought it, who asked it — comes afterwards, so the reader hears the words first and learns whose they were second.

Two things follow from that, and both are worth knowing before you build a sentence with it:

1. \textbf{The quoted words are not changed.} English shifts tense and person when it reports — \emph{he said he was going}. Sanskrit does not. It hands you the words as they were spoken and lets \textbf{\sk{इति}} mark where they stop. 2. \textbf{\sk{इति} is not a verb and never inflects.} It is the same three letters in every sentence you will ever meet it in.

You have probably met it already without knowing. \textbf{\sk{इत्यादि}} (\emph{ityādi}) — \emph{\textquotedblleft{}and so on\textquotedblright{}} — is \textbf{\sk{इति}} plus \textbf{\sk{आदि}}, \textquotedblleft{}beginning\textquotedblright{}: \emph{thus, and the rest}. And the closing line of every chapter of a classical text begins \textbf{\sk{इति}} — \emph{thus ends…}.

\begin{grammarlens}[title={What you've built}]
The mark that closes a quotation, in a language with no quotation marks.
\end{grammarlens}

\subsection*{Your turn}
\begin{itemize}
\raggedright
  \item \emph{Say it:} \emph{iti}
  \item \emph{Read it:} \textbf{\sk{इति}}, and say which side of it the quoted words sit on
  \item \emph{Recall:} what \textbf{\sk{इत्यादि}} is made of
\end{itemize}

\begin{tcolorbox}[breakable,colback=teal!4,colframe=teal!35!black,title={Before you move on}]
Where do the quoted words stand relative to \textbf{\sk{इति}}? (Before it.) Say \textbf{\sk{इति}} once more.
\end{tcolorbox}
\section[… iti cintayati]{… \sk{इति} \sk{चिन्तयति} — a clause as the object}

\label{lesson:SA-C54-iti-clause}



\begin{quote}
Say \textbf{\sk{इति}}, and say which side of it the quoted words sit on.
\end{quote}

\begin{grammarlens}[title={Grammar Lens: what a thinking verb thinks}]
\textbf{\sk{चिन्तयति}} has been in this book since the thinking verb arrived, with nothing to think. Every verb you have used takes a \emph{noun} as its object — you see a house, you eat food. \textbf{\sk{इति}} lets a verb take a \textbf{whole sentence} instead.

\begin{quote}
\textbf{\sk{बालः} \sk{गच्छति} \sk{इति} \sk{गुरुः} \sk{चिन्तयति}}
\end{quote}

\begin{quote}
\emph{bālaḥ gacchati iti guruḥ cintayati}
\end{quote}

\begin{quote}
\textbf{the teacher thinks that the boy is going}
\end{quote}

Take it apart in the order Sanskrit gives it to you:

\noindent
\begin{tabularx}{\linewidth}{@{}>{\raggedright\arraybackslash}X>{\raggedright\arraybackslash}X>{\raggedright\arraybackslash}X@{}}
\toprule
\textbf{piece} & \textbf{meaning} & \textbf{job} \\
\midrule
\textbf{\sk{बालः} \sk{गच्छति}} & the boy goes & the clause being thought \\
\textbf{\sk{इति}} & thus & the clause ends here \\
\textbf{\sk{गुरुः} \sk{चिन्तयति}} & the teacher thinks & the frame around it \\
\bottomrule
\end{tabularx}

The inner clause is a complete sentence and stays one. Nothing is bent to make it fit inside the outer one — no \emph{that}, no change of ending, no reordering. \textbf{\sk{इति}} does all the joining by standing between them.

Swap the outer verb and you have every reporting sentence at once: \textbf{… \sk{इति} \sk{वदति}} (says), \textbf{… \sk{इति} \sk{पृच्छति}} (asks), \textbf{… \sk{इति} \sk{जानाति}} (knows). The inner clause never changes.

Compare this with the chapter before. \textbf{\sk{यः} … \sk{सः}} joined two clauses by putting a pronoun in each. \textbf{\sk{इति}} joins two clauses by putting one particle at the seam. Two joins, two mechanisms — and neither of them is a conjunction in the English sense.
\end{grammarlens}

\begin{grammarlens}[title={What you've built}]
A sentence that reports another sentence.
\end{grammarlens}

\subsection*{Your turn}
\begin{itemize}
\raggedright
  \item \emph{Say it:} \emph{bālaḥ gacchati iti guruḥ cintayati}
  \item \emph{Say it:} it again, and name where the inner clause stops
  \item \emph{Read it:} \textbf{\sk{बालः} \sk{गच्छति} \sk{इति} \sk{गुरुः} \sk{चिन्तयति}}, then say what it means
  \item \emph{Contrast:} the correlative puts a pronoun in each clause; \textbf{\sk{इति}} puts one particle at the seam
\end{itemize}

\begin{tcolorbox}[breakable,colback=teal!4,colframe=teal!35!black,title={Before you move on}]
Does the inner clause change when it is reported? (No.) Say the teacher-and-boy sentence once more.
\end{tcolorbox}
\section[iti]{\sk{इति} — reviewed}

\label{lesson:SA-R54-iti}



\begin{quote}
Say \textbf{\sk{इति}}, and say what it does to the words in front of it.
\end{quote}

\begin{grammarlens}[title={What you've built}]
Two chapters, two ways of putting two clauses in one sentence, and they work differently on purpose.

The correlative marks \textbf{both} clauses — a \textbf{\sk{य}-} word in one, a \textbf{\sk{त}-} word in the other — and the clauses stay equals. \textbf{\sk{इति}} marks \textbf{one} seam, and makes the clause in front of it into the object of the verb behind it.

That is the difference between \emph{which one goes, that one sees} and \emph{the teacher thinks that the boy goes}. In the first, neither clause is inside the other. In the second, the whole first clause is what the second one is about.
\end{grammarlens}

\subsection*{Your turn}
\begin{itemize}
\raggedright
  \item \emph{Say it:} \emph{bālaḥ gacchati iti guruḥ cintayati}
  \item \emph{Say it:} \emph{yaḥ gacchati saḥ paśyati}
  \item \emph{Contrast:} which sentence puts one clause INSIDE the other
\end{itemize}

\begin{tcolorbox}[breakable,colback=teal!4,colframe=teal!35!black,title={Before you move on}]
Name the two joining mechanisms you now have, and say which one subordinates.
\end{tcolorbox}
